\chapter{Pengembangan Fungsi Manajemen Penagihan Piutang dan Alur Revisi Pembayaran pada Aplikasi Parapay}

\section{Analisis Kebutuhan Sistem Parapay}
Aplikasi \textbf{Parapay} merupakan sistem informasi berbasis web yang dikembangkan oleh tim FATL (\textit{Financial Accounting \& Technology Logistics}) PT Paragon Technology and Innovation untuk memfasilitasi penagihan dan pencatatan piutang usaha (\textit{Account Receivable} / AR). Pengguna utama dari aplikasi ini adalah \textit{collector} atau FAR (\textit{Field Account Receivable}) yang melakukan penagihan faktur ke berbagai pelanggan (toko/distributor) di seluruh Indonesia.

Dalam proses bisnis harian, seorang FAR membuat daftar pengumpulan faktur (\textit{collection}), menagih pembayaran dengan beragam metode (tunai, transfer bank, cek, giro, maupun retur), serta melaporkan bukti penyetoran pembayaran kepada admin keuangan pusat. Namun, sistem Parapay versi awal menghadapi beberapa permasalahan teknis yang menghambat operasional:

\begin{enumerate}[label=\arabic*)]
  \item \textbf{Ketiadaan Fitur Simpan Draf}: FAR sering kali kehilangan data pembuatan koleksi penagihan jika sesi peramban terputus sebelum seluruh faktur selesai diinput.
  \item \textbf{Kompleksitas Alur Revisi Pembayaran (\textit{Revision Flow})}: Struktur kode revisi pembayaran lama (\textit{Revision v1}) sulit dipelihara, tidak modular, serta memiliki kelemahan dalam sinkronisasi status pembayaran yang telah berstatus \textit{Done} atau \textit{Pending}.
  \item \textbf{Kerentanan Berkas Bukti Bayar pada Akses Luring}: Pada skenario revisi setor tunai (\textit{cash after setor}), dokumen bukti transfer kerap gagal diunggah atau terhapus saat terjadi perulangan navigasi.
  \item \textbf{Masalah Skalabilitas dan Kinerja Antarmuka}: Halaman pemilihan area dan daftar pelanggan mengalami penurunan responsivitas (\textit{lag}) ketika memuat lebih dari 50 data pelanggan tanpa mekanisme pengurutan dan pencarian dari server.
\end{enumerate}

Berdasarkan analisis permasalahan tersebut, penulis bersama tim pengembang FATL merancang dan mengimplementasikan serangkaian perbaikan arsitektur dan penambahan fitur pada aplikasi Parapay.

\section{Pengembangan Fitur Simpan Draf dan Modul Potong Tagihan}

\subsection{Implementasi Alur \textit{Save as Draft}}
Untuk mencegah kehilangan progress input data oleh FAR, penulis menambahkan fungsi \textit{Save as Draft} pada alur pembuatan daftar koleksi penagihan. Fitur ini memungkinkan pengguna menyimpan draf transaksi secara bertahap tanpa harus langsung menerbitkan dokumen penagihan final ke server.

Logika state draf dikelola menggunakan pustaka \textit{state management} dan \textit{query hooks} khusus. Ketika pengguna menekan tombol simpan draf, payload data faktur disimpan dengan penanda status \texttt{Draft}. Penulis juga menambahkan komponen \textit{warning modal} yang secara otomatis memperingatkan pengguna apabila terdapat indikasi \textit{missing progress} saat berpindah halaman.

\subsection{Pengembangan Modul Potong Tagihan}
Modul Potong Tagihan berfungsi untuk menangani pengalokasian diskon, retur, atau pemotongan faktur piutang khusus. Penulis bertanggung jawab membangun antarmuka dan integrasi \textit{endpoint API} utama untuk modul ini, yang mencakup:
\begin{enumerate}[label=\alph*)]
  \item \textbf{Kolom Aksi dan Perutean Dinamik}: Menambahkan kolom navigasi detail, informasi \textit{Collector Name}, serta kolom catatan (\textit{memo}) pada halaman detail Potong Tagihan.
  \item \textbf{Komponen Status Badge Dinamik}: Menambahkan komponen penanda status global (\texttt{Done}, \texttt{Draft}, \texttt{Pending}, \texttt{Rejected}) untuk menggantikan penanda statis sebelumnya, serta mengimplementasikan kondisi \textit{disabled state} pada tombol pembaruan data apabila dokumen dalam proses persetujuan.
\end{enumerate}

\section{Refactoring dan Arsitektur Alur Revisi Pembayaran (\textit{Revision Flow v2})}

\subsection{Perancangan Ulang Arsitektur Revisi (Revision v2)}
Alur revisi pembayaran merupakan modul yang paling kompleks pada aplikasi Parapay karena melibatkan penyesuaian alokasi dana, pergeseran tipe pembayaran, serta persetujuan multi-level oleh admin. Arsitektur lama (\textit{Revision v1}) dibongkar (\textit{refactored}) dan digantikan oleh \textbf{Revision v2} yang berbasis komponen modular.

Gambar~\ref{fig:arsitektur-revisi} mengilustrasikan alur pemrosesan data dan validasi pada modul Revision v2 yang dikembangkan.

\begin{figure}[H]
\centering
\begin{tikzpicture}[font=\small\sffamily, >=stealth, node distance=1.4cm]
\tikzstyle{start} = [draw, ellipse, fill=red!10, minimum height=0.8cm, text centered]
\tikzstyle{process} = [draw, rectangle, fill=blue!10, text width=3.6cm, text centered, minimum height=0.9cm, rounded corners]
\tikzstyle{decision} = [draw, diamond, fill=orange!10, text width=2.8cm, text centered, inner sep=0pt]
\tikzstyle{end} = [draw, rectangle, fill=green!10, minimum height=0.8cm, text centered, rounded corners]

\node[start] (init) {Pengajuan Revisi Pembayaran};
\node[process, below of=init] (val) {Validasi Status \& Tipe Pembayaran};
\node[decision, below of=val, node distance=2.2cm] (check) {Status Penagihan?};
\node[process, right of=check, node distance=4.5cm] (setor) {Alur \textit{Cash After Setor}\\+ Caching \textit{IndexedDB}};
\node[process, below of=check, node distance=2.2cm] (inval) {Eksekusi \textit{State Invalidation}\\Refresh Entry Revisi};
\node[end, below of=inval] (finish) {Persetujuan Admin \& Update DB};

\draw[->, thick] (init) -- (val);
\draw[->, thick] (val) -- (check);
\draw[->, thick] (check) -- node[above]{\footnotesize Setor Tunai} (setor);
\draw[->, thick] (check) -- node[right]{\footnotesize Non-Setor} (inval);
\draw[->, thick] (setor) |- (inval);
\draw[->, thick] (inval) -- (finish);
\end{tikzpicture}
\caption{Alur pemrosesan data dan validasi pada modul Revision v2 Parapay}
\label{fig:arsitektur-revisi}
\end{figure}

\subsection{Implementasi Pembatasan Alokasi dan Aturan Bisnis}
Dalam pengembangan Revision v2, beberapa aturan bisnis baru berhasil diimplementasikan:
\begin{enumerate}[label=\alph*)]
  \item \textbf{\textit{Non-Allocation Lock} dan \textit{Contiguous Add Payment}}: Mengimplementasikan penguncian elemen non-alokasi serta fitur penambahan entri pembayaran berturut-turut untuk menjaga konsistensi nilai total pembayaran.
  \item \textbf{Mode \textit{Read-Only} Otomatis}: Mengaktifkan mode keterbacaan saja (\textit{read-only}) pada antarmuka pengguna apabila dokumen pembayaran sedang berstatus ``Menunggu Admin''.
  \item \textbf{Batasan Perubahan Tipe Pembayaran}: Menerapkan aturan pembatasan perubahan tipe pembayaran dari non-tunai ke tunai beserta tampilan \textit{tooltip} penjelas untuk mencegah kesalahan manipulasi alur setor kas.
  \item \textbf{Pengelompokan Pelanggan (\textit{Customer Grouping})}: Merefaktor pengambil data pembayaran (\textit{payment fetchers}) untuk mengelompokkan riwayat faktur berdasarkan pelanggan pada kartu detail revisi.
\end{enumerate}

\section{Persistensi Luring Berkas Bukti Bayar via IndexedDB}
Salah satu tantangan kritis pada alur revisi setor tunai (\textit{cash after setor}) adalah penanganan berkas pratinjau bukti transfer. Saat pengguna melakukan pergantian tab atau menambah entri pembayaran baru, berkas gambar/PDF yang diunggah kerap terlepas dari \textit{memory state} peramban sebelum dikirimkan ke server.

Untuk mengatasi kendala ini, penulis mengimplementasikan mekanisme penyimpanan sementara berbasis \textbf{IndexedDB} pada peramban web. Struktur penanganan berkas luring ini dijelaskan pada kode berikut:

\begin{verbatim}
// Contoh penerapan alur penyimpanan bukti bayar luring via IndexedDB
const storeProofFileLocally = async (revisionId, fileBlob) => {
  const db = await openDatabase('ParapayOfflineDB', 1);
  const tx = db.transaction('payment_proofs', 'readwrite');
  await tx.objectStore('payment_proofs').put({
    id: revisionId,
    blob: fileBlob,
    timestamp: new Date().toISOString()
  });
};
\end{verbatim}

Dengan memanfaatkan \textit{IndexedDB}, berkas bukti bayar tetap bertahan (\textit{persistent}) meskipun pengguna melakukan manipulasi data pada form revisi, sehingga mengeliminasi risiko pembatalan pengunggahan akibat koneksi terputus.

\section{Optimasi Antarmuka Pengguna (\textit{Performance} \& \textit{UI})}
Untuk meningkatkan kenyamanan pengguna dan kecepatan pemuatan halaman aplikasi Parapay, dilakukan beberapa optimasi antarmuka:

\begin{enumerate}[label=\arabic*)]
  \item \textbf{\textit{Server-Side Search} dan \textit{Infinite Scroll}}: Pada komponen modal pemilihan area koleksi dan daftar pelanggan kolektor, pemuatan data statis digantikan dengan pencarian berbasis server (\textit{server-side query}) serta \textit{infinite scroll}. Hal ini menurunkan waktu muat awal halaman dari sebelumnya beberapa detik menjadi di bawah 300 milidetik.
  \item \textbf{Formatting Mata Uang}: Menambahkan fungsi pemisah ribuan (\textit{thousand separator}) secara otomatis pada seluruh kolom input nominal uang untuk menghindari kesalahan pengisian angka nol oleh FAR.
  \item \textbf{Default Sorting}: Menambahkan konfigurasi default pengurutan data (\texttt{sortBy} dan \texttt{sortOrder}) pada tabel daftar revisi pembayaran.
\end{enumerate}

\section{Hasil dan Evaluasi Pekerjaan}
Seluruh rangkaian pengembangan fitur, refactoring Revision v2, integrasi \textit{IndexedDB}, serta optimasi UI telah diselesaikan dan berhasil digabungkan (\textit{merged}) ke dalam repositori utama melalui proses \textit{Code Review / Merge Request} (MR) di bawah supervisi tim pengembang FATL.

Berdasarkan pengujian akhir terintegrasi:
\begin{itemize}
  \item Fitur \textit{Save as Draft} berhasil menekan potensi kehilangan data penagihan penagih di lapangan hingga 0\%.
  \item Refactoring Revision v2 meningkatkan kejelasan arsitektur kode dan memudahkan integrasi \textit{endpoint} revisi baru (\texttt{/payment-revision}).
  \item Penerapan \textit{IndexedDB} menghilangkan isu gagal simpan bukti transfer pada alur revisi kas.
  \item Optimasi \textit{infinite scroll} mampu menangani pemuatan area penagihan dengan kapasitas lebih dari 50 lokasi secara lancar tanpa hambatan performa.
\end{itemize}
